\chapter{Desarrollo del ETL}

\section{Extracción}

La extracción se coordina desde \texttt{src/main.py} y \texttt{src/extract.py}. El proceso usa \texttt{data/raw/metadata\_fuentes.csv} como contrato operativo para ubicar archivo, hoja, \texttt{skiprows}, formato y observaciones por fuente. Las fuentes A y B se leen mediante un parser específico para SpreadsheetML/XML 2003, mientras que las fuentes C y D se leen con \texttt{pandas.read\_excel()} respetando la configuración registrada en metadata.

El resultado de esta etapa no es solo la lectura en memoria, sino también la generación de evidencia diagnóstica en \texttt{outputs/perfilado\_fuentes.xlsx} y \texttt{outputs/conflictos\_fuentes.md}. Esto permite detectar estructura real, filas fuera de alcance, valores especiales y heterogeneidad en nombres comunales antes de construir los staging.

\section{Transformación}

La transformación se implementa en \texttt{src/transform.py}. Cada fuente produce un staging independiente y validado. Las reglas aplicadas incluyen selección de columnas, renombre semántico, tipificación, formateo de \texttt{codigo\_comuna}, estandarización de \texttt{nombre\_comuna}, tratamiento de valores especiales, filtrado al universo final de 32 comunas y validaciones de rango.

\begin{table}[htbp]
\centering
\footnotesize
\caption{Transformaciones y controles relevantes observados en el ETL.}
\begin{tabular}{|p{2.7cm}|p{1.2cm}|p{7.0cm}|p{3.5cm}|}
\hline
Transformación & Fuente & Regla aplicada & Evidencia \\
\hline
Selección de columnas & A-D & Solo se conservan campos comunales y variables de interés. & outputs/resumen\_transformaciones.md \\
Normalización comunal & A-D & Estandarización de nombre\_comuna contra dim\_comuna\_base.csv. & outputs/resumen\_transformaciones.md; outputs/homologacion\_comunas.csv \\
Formateo de llave & A-D & codigo\_comuna se tipifica y luego se fija como string de 5 dígitos. & src/transform.py; src/integrate.py \\
Tratamiento de valores especiales & A & No Aplica -> 0 y No Recepcionado -> NA en componentes de áreas verdes. & outputs/conflictos\_fuentes.md; src/transform.py \\
Conversión de escala & C & La pobreza se transforma desde proporción 0-1 a porcentaje 0-100. & data/raw/metadata\_fuentes.csv; src/transform.py \\
Filtrado territorial & A-D & Se excluyen comunas fuera de la Provincia de Santiago y filas sin código válido. & outputs/resumen\_transformaciones.md \\
Renombre semántico & B-D & iadm41\_2024 -> ipp\_miles\_pesos; Población censada -> poblacion. & src/transform.py \\
Métricas derivadas & Fase 6 & Se calculan areas\_verdes\_m2\_hab e ipp\_pesos\_hab con control de población > 0. & outputs/resumen\_dataset\_final.md; src/integrate.py \\
Validaciones de rango & A-D & Población > 0; pobreza 0-100; áreas verdes e IPP no negativos. & outputs/validacion\_staging.csv; outputs/reporte\_validacion\_final.csv \\
Control de cardinalidad & Fase 6 & Todos los merges se ejecutan como left con validación one\_to\_one. & outputs/log\_integracion.md; src/integrate.py \\
\hline
\end{tabular}
\end{table}


\section{Integración}

La integración se implementa en \texttt{src/integrate.py}. El proceso parte desde \texttt{dim\_comuna\_base.csv} y ejecuta cuatro merges \texttt{left} con validación \texttt{one\_to\_one}, en el orden población, pobreza, áreas verdes y capacidad municipal. Después del merge se asignan los años de referencia por variable y se calculan las métricas derivadas \texttt{areas\_verdes\_m2\_hab} e \texttt{ipp\_pesos\_hab}, controlando divisiones inválidas e infinitos.

El flujo completo puede reejecutarse desde la raíz del repositorio con \texttt{python src/main.py} o \texttt{python -m src.main}. Los reportes \texttt{outputs/log\_integracion.md}, \texttt{outputs/resumen\_dataset\_final.md} y \texttt{outputs/reporte\_correcciones\_pre\_fase10.md} respaldan tanto la lógica del merge como la reproducibilidad observada del pipeline.
